\section{الحسابات الديناميكية وضبط المعايير \textenglish{(Kinematics \& Parameterization)}}
\label{sec:kinematics_and_parameters}

يُعد هذا المبحث العصب الرياضي لمنظومة المقص الطائر، حيث يتم تحويل المتطلبات الإنتاجية (سرعة السير وطول المنتج) إلى قيم نبضية وترددات رقمية يتعامل معها جهاز الـ \textenglish{PLC}. تعتمد الحسابات على تحقيق التزامن التام بين المحور القائد (المشفر) والمحور التابع (محرك السيرفو).

\subsection{حسابات المحور القائد \textenglish{(Master Axis - Encoder)}}
يعتمد النظام على قراءة سرعة وموقع سير الصابون عبر المشفر النبضي لتحديد لحظة القص.
\begin{itemize}
    \item \textbf{دقة المشفر الميكانيكية:} تمت معايرة المشفر بحيث يُعطي $1000\text{ pulses}$ (نبضة) لكل مسافة مقطوعة قدرها $100\text{ mm}$ من حركة السير الناقل.
    \item \textbf{الطول المطلوب للمنتج:} $20\text{ cm}$ وهو ما يعادل $200\text{ mm}$.
    \item \textbf{حساب نبضات القص:} لمعرفة القيمة التي يجب أن يصل إليها العداد السريع (\textenglish{C254}) لاتخاذ قرار القص، نستخدم التناسب الخطي:
    $$x = \frac{200\text{ mm} \times 1000\text{ pulses}}{100\text{ mm}} = 2000\text{ pulses}$$
    إذن، الشرط البرمجي لتفعيل عملية القص هو وصول العداد إلى $2000$ نبضة.
\end{itemize}

\subsection{حسابات المحور التابع \textenglish{(Slave Axis - Servo Motor)}}
يُستخدم محرك سيرفو من طراز \textenglish{ASDA-B2} باستطاعة $400\text{ W}$ لتحريك عربة المقص.
\begin{itemize}
    \item \textbf{دقة المحرك الافتراضية:} يحتاج المحرك إلى $10,000\text{ pulses}$ لإتمام دورة ميكانيكية كاملة (\textenglish{1 Rev}).
    \item \textbf{الخطوة الميكانيكية (\textenglish{Pitch/Lead}):} كل دورة كاملة للمحرك تُحرك العربة مسافة $10\text{ mm}$.
    \item \textbf{النبضات لكل مليمتر:} 
    $$\text{Pulses per mm} = \frac{10,000\text{ pulses}}{10\text{ mm}} = 1000\text{ pulses/mm}$$
\end{itemize}

\subsection{سرعة التزامن والتردد المطلوب}
لتحقيق التزامن الحركي التام أثناء القص، يجب أن تكون السرعة الخطية لعربة المقص مساوية تماماً للسرعة الخطية لسير الصابون، والتي تبلغ في هذه المنظومة $\mathbf{150\text{ mm/s}}$.
\begin{itemize}
    \item \textbf{حساب تردد المزامنة:} لمعرفة التردد النبضي الذي يجب أن يرسله الـ \textenglish{PLC} للمحرك لكي ينطلق بهذه السرعة، نقوم بضرب السرعة الخطية في عدد النبضات اللازمة لقطع مليمتر واحد:
    $$f = v \times \text{Pulses per mm} = 150\text{ mm/s} \times 1000\text{ pulses/mm} = 150,000\text{ Hz}$$
\end{itemize}

التردد الناتج ($150\text{ kHz}$) يقع ضمن النطاق التشغيلي الآمن لمخارج الترانزستور في متحكم \textenglish{Delta DVP-28SV11T} (والذي يدعم حتى $200\text{ kHz}$). بناءً على ذلك، تم اعتماد هذا التردد في تعليمة توليد النبضات (\LR{DPLSY}) لضمان انطلاق السيرفو بسرعة التطابق المطلوبة بشكل سلس ومستقر.

\subsection{التحليل الهندسي لترددات التشغيل \textenglish{(Operating Frequencies Analysis)}}
يعكس هذا التصميم توازناً هندسياً دقيقاً ومدروساً بين متطلبات القراءة (المداخل) ومتطلبات الحركة (المخارج) في جهاز الـ \textenglish{PLC}، حيث تم تجنب رفع الترددات بشكل عشوائي:

\begin{itemize}
    \item \textbf{تردد القراءة الآمن (المحور القائد):} بضبط دقة المشفر على $10\text{ pulses/mm}$ (أي $1000$ نبضة لكل $100\text{ mm}$) وسرعة السير $150\text{ mm/s}$، فإن التردد الداخل للعداد السريع يبلغ $1500\text{ Hz}$ ($1.5\text{ kHz}$) فقط. هذا التردد المنخفض نسبياً يضمن استقراراً عالياً في القراءة ومناعة ممتازة ضد التشويش الكهرومغناطيسي \textenglish{(EMI)} في البيئة الصناعية، مع الاحتفاظ بدقة قياس ممتازة للتطبيق ($0.1\text{ mm}$).
    
    \item \textbf{تردد التحكم الدقيق (المحور التابع):} على الجانب الآخر، ولضمان أقصى درجات النعومة والدقة في حركة السيرفو أثناء القص، تم إبقاء دقة المحرك عالية جداً عند $1000\text{ pulses/mm}$. لتحقيق سرعة المزامنة ($150\text{ mm/s}$)، يقوم الـ \textenglish{PLC} بتوليد نبضات بتردد $150,000\text{ Hz}$ ($150\text{ kHz}$).
\end{itemize}

تُثبت هذه الأرقام كفاءة توظيف متحكم \textenglish{Delta DVP-28SV11T}، والذي استطاع الربط بين إشارة قراءة منخفضة ومستقرة من جهة، وتوليد إشارة تحكم سريعة وعالية الدقة من جهة أخرى، وكلاهما يعملان بشكل متزامن وبأمان تام تحت سقف القدرة العتادية القصوى للمنافذ ($200\text{ kHz}$).